\documentclass[12pt,a4paper]{article}

\usepackage{fontspec}
\usepackage{polyglossia}
\setmainfont{Liberation Serif}
\setsansfont{Liberation Sans}
\setmonofont{Liberation Mono}
\setdefaultlanguage{russian}
\usepackage{amsmath,amssymb}
\usepackage{geometry}
\usepackage{booktabs}

\geometry{margin=2.2cm}
\setlength{\parskip}{0.35em}
\setlength{\parindent}{1.25em}

\title{Автономная система двух уравнений.\\
Вариант~13: особые точки и метод Рунге--Кутты 4-го порядка}
\author{}
\date{}

\begin{document}
\maketitle

\section{Постановка}

Рассматривается автономная система
\begin{equation}
\label{eq:system}
\begin{cases}
\dfrac{dx}{dt} = f(x,y) = (2x-y)^2 - 9,\\[1mm]
\dfrac{dy}{dt} = g(x,y) = 9 - (x-2y)^2.
\end{cases}
\end{equation}
Правая часть не зависит от \(t\) явно, поэтому качественная картина полностью определяется векторным полем \((f,g)\) на плоскости \((x,y)\).

Задача состоит из двух частей.
\begin{enumerate}
  \item Найти особые точки системы~\eqref{eq:system}, линеаризовать поле в их окрестности и классифицировать точки по собственным значениям матрицы Якоби.
  \item Проинтегрировать траектории численно методом~2 четвёртого порядка (правило трёх восьмых). Шаг по методичке \(h = 0{,}05\). Локальная погрешность одного шага имеет порядок \(O(h^5)\), погрешность на конечном отрезке~--- порядок \(O(h^4)\).
\end{enumerate}

\section{Особые точки}

Особая точка \((x_*, y_*)\)~--- положение равновесия:
\[
f(x_*, y_*) = 0, \qquad g(x_*, y_*) = 0.
\]
Это равносильно системе
\begin{equation}
\label{eq:alg}
(2x-y)^2 = 9, \qquad (x-2y)^2 = 9,
\end{equation}
то есть
\[
2x - y = \pm 3, \qquad x - 2y = \pm 3.
\]
Удобна замена
\begin{equation}
\label{eq:uv}
u = 2x - y, \qquad v = x - 2y.
\end{equation}
Она обратима:
\begin{equation}
\label{eq:xy}
x = \frac{2u - v}{3}, \qquad y = \frac{u - 2v}{3},
\end{equation}
поскольку \(2u - v = 3x\) и \(u - 2v = 3y\). Условия~\eqref{eq:alg} принимают вид \(u = \pm 3\), \(v = \pm 3\). Четыре комбинации знаков дают все положения равновесия.

\begin{center}
\begin{tabular}{cccc}
\toprule
\(u\) & \(v\) & \(x = (2u-v)/3\) & \(y = (u-2v)/3\) \\
\midrule
\(3\)  & \(3\)  & \(1\)  & \(-1\) \\
\(3\)  & \(-3\) & \(3\)  & \(3\) \\
\(-3\) & \(3\)  & \(-3\) & \(-3\) \\
\(-3\) & \(-3\) & \(-1\) & \(1\) \\
\bottomrule
\end{tabular}
\end{center}

Других вещественных решений нет: каждое из уравнений~\eqref{eq:alg} задаёт пару прямых, и прямые первой пары пересекают прямые второй ровно в этих четырёх точках. Итак, особые точки
\begin{equation}
\label{eq:points}
(1,-1),\quad (3,3),\quad (-3,-3),\quad (-1,1).
\end{equation}

\section{Линеаризация}

Поведение траекторий вблизи изолированной особой точки совпадает с поведением линейной системы
\[
\begin{pmatrix} \xi' \\ \eta' \end{pmatrix}
=
J(x_*, y_*)
\begin{pmatrix} \xi \\ \eta \end{pmatrix},
\qquad
\xi = x - x_*,\quad \eta = y - y_*,
\]
если \(\det J(x_*, y_*) \neq 0\). Матрица Якоби поля \((f,g)\):
\[
J(x,y) =
\begin{pmatrix}
f_x & f_y \\
g_x & g_y
\end{pmatrix}.
\]
Дифференцируем~\eqref{eq:system}:
\begin{align*}
f_x &= 2(2x-y)\cdot 2 = 4(2x-y), &
f_y &= 2(2x-y)\cdot(-1) = -2(2x-y),\\
g_x &= -2(x-2y)\cdot 1 = -2(x-2y), &
g_y &= -2(x-2y)\cdot(-2) = 4(x-2y).
\end{align*}
Следовательно,
\begin{equation}
\label{eq:J}
J(x,y) =
\begin{pmatrix}
4(2x-y) & -2(2x-y) \\
-2(x-2y) & 4(x-2y)
\end{pmatrix}.
\end{equation}
В особой точке \(u = 2x-y = \pm 3\) и \(v = x-2y = \pm 3\), поэтому матрица невырождена во всех четырёх точках, и классификация по линейному приближению законна.

Тип точки определяется корнями характеристического уравнения
\begin{equation}
\label{eq:char}
\lambda^2 - \tau\lambda + \Delta = 0,
\qquad
\tau = \operatorname{tr} J,\quad
\Delta = \det J.
\end{equation}
Дискриминант \(D = \tau^2 - 4\Delta\).
\begin{itemize}
  \item \(\Delta < 0\): собственные значения вещественны и разных знаков~--- седло.
  \item \(\Delta > 0\), \(D > 0\), \(\tau > 0\): неустойчивый узел.
  \item \(\Delta > 0\), \(D > 0\), \(\tau < 0\): устойчивый узел.
  \item \(\Delta > 0\), \(D < 0\), \(\tau \neq 0\): фокус (устойчивый при \(\tau < 0\)).
  \item \(\Delta > 0\), \(\tau = 0\): центр линейного приближения.
\end{itemize}
В данной задаче комплексных собственных значений нет: все четыре точки~--- узлы или сёдла.

\section{Классификация}

\subsection{Точка \((1,-1)\): неустойчивый узел}

Здесь \(u = 3\), \(v = 3\), поэтому
\[
J(1,-1) =
\begin{pmatrix} 12 & -6 \\ -6 & 12 \end{pmatrix}.
\]
След и определитель:
\[
\tau = 24, \qquad \Delta = 144 - 36 = 108, \qquad D = 576 - 432 = 144 > 0.
\]
Оба корня положительны:
\[
\lambda = \frac{24 \pm 12}{2},
\qquad
\lambda_1 = 18, \quad \lambda_2 = 6.
\]
Проверка: \(\lambda^2 - 24\lambda + 108 = (\lambda-18)(\lambda-6)\).

Собственные векторы.
\begin{align*}
(J - 18I)\,v &= 0
&&\Longrightarrow&&
v_1 + v_2 = 0,
&& v^{(1)} = (1,-1),\\
(J - 6I)\,v &= 0
&&\Longrightarrow&&
v_1 = v_2,
&& v^{(2)} = (1,1).
\end{align*}
Траектории уходят от точки по обоим направлениям. Вдоль \((1,-1)\) уход быстрее: показатель \(e^{18t}\) растёт быстрее, чем \(e^{6t}\). Это неустойчивый узел.

\subsection{Точка \((-1,1)\): устойчивый узел}

Здесь \(u = v = -3\) и
\[
J(-1,1) =
\begin{pmatrix} -12 & 6 \\ 6 & -12 \end{pmatrix}
= -J(1,-1).
\]
Значит,
\[
\tau = -24, \qquad \Delta = 108, \qquad \lambda_1 = -6, \quad \lambda_2 = -18.
\]
Оба корня отрицательны, точка асимптотически устойчива. Направление быстрого приближения~--- \((1,-1)\) (показатель \(e^{-18t}\)), направление медленного~--- \((1,1)\) (показатель \(e^{-6t}\)). Почти все траектории входят в точку, касаясь более медленного направления \((1,1)\).

\subsection{Точка \((3,3)\): седло}

Здесь \(u = 3\), \(v = -3\) и
\[
J(3,3) =
\begin{pmatrix} 12 & -6 \\ 6 & -12 \end{pmatrix}.
\]
\[
\tau = 0, \qquad \Delta = -144 + 36 = -108 < 0.
\]
Характеристическое уравнение \(\lambda^2 - 108 = 0\):
\[
\lambda_{1,2} = \pm\sqrt{108} = \pm 6\sqrt{3} \approx \pm 10{,}3923.
\]
Собственные направления.
\begin{align*}
\lambda = 6\sqrt{3} > 0:&\quad
v^{(u)} = \bigl(1,\; 2-\sqrt{3}\bigr) \approx (1;\; 0{,}268),\\
\lambda = -6\sqrt{3} < 0:&\quad
v^{(s)} = \bigl(1,\; 2+\sqrt{3}\bigr) \approx (1;\; 3{,}732).
\end{align*}
Вдоль \(v^{(u)}\) траектории уходят, вдоль \(v^{(s)}\)~--- приходят. Это седло. Устойчиво оно только на одномерной сепаратрисе, поэтому для задачи Коши с произвольными начальными данными точка неустойчива.

\subsection{Точка \((-3,-3)\): седло}

Здесь \(u = -3\), \(v = 3\) и
\[
J(-3,-3) =
\begin{pmatrix} -12 & 6 \\ -6 & 12 \end{pmatrix}.
\]
Снова \(\tau = 0\), \(\Delta = -108\), собственные значения \(\pm 6\sqrt{3}\). Направления меняются местами относительно предыдущего седла:
\begin{align*}
\lambda = 6\sqrt{3} > 0:&\quad
v^{(u)} = \bigl(1,\; 2+\sqrt{3}\bigr),\\
\lambda = -6\sqrt{3} < 0:&\quad
v^{(s)} = \bigl(1,\; 2-\sqrt{3}\bigr).
\end{align*}
Точка~--- седло.

\subsection{Сводка}

\begin{center}
\begin{tabular}{cccc}
\toprule
Точка & \(\tau\) & \(\Delta\) & Тип \\
\midrule
\((1,-1)\)   & \(24\)  & \(108\)  & неустойчивый узел, \(\lambda = 18\) и \(6\) \\
\((3,3)\)    & \(0\)   & \(-108\) & седло, \(\lambda = \pm 6\sqrt{3}\) \\
\((-3,-3)\)  & \(0\)   & \(-108\) & седло, \(\lambda = \pm 6\sqrt{3}\) \\
\((-1,1)\)   & \(-24\) & \(108\)  & устойчивый узел, \(\lambda = -6\) и \(-18\) \\
\bottomrule
\end{tabular}
\end{center}

Единственная асимптотически устойчивая особая точка~--- \((-1,1)\). Из окрестности \((1,-1)\) решения уходят. Вблизи сёдел сохраняются только две сепаратрисы.

\section{Метод 2 четвёртого порядка}

В методичке метод записан для одного уравнения \(y' = f(x,y)\), где \(x\)~--- независимая переменная. У системы~\eqref{eq:system} неизвестны обе координаты, а время в правую часть не входит. Поэтому узлы \(h/3\), \(2h/3\), \(h\) задают только момент стадии, а сдвигаются обе компоненты вектора состояния.

Обозначим \(F = (f,g)\) и \(Q_i = (Q_i^{x}, Q_i^{y})\). Правило трёх восьмых, оно же метод~2 при весе \((1,3,3,1)/8\):
\begin{align}
Q_1 &= h\, F(x_n, y_n), \label{eq:q1}\\
Q_2 &= h\, F\Bigl(x_n + \tfrac13 Q_1^{x},\; y_n + \tfrac13 Q_1^{y}\Bigr), \label{eq:q2}\\
Q_3 &= h\, F\Bigl(x_n - \tfrac13 Q_1^{x} + Q_2^{x},\;
                  y_n - \tfrac13 Q_1^{y} + Q_2^{y}\Bigr), \label{eq:q3}\\
Q_4 &= h\, F\bigl(x_n + Q_1^{x} - Q_2^{x} + Q_3^{x},\;
                  y_n + Q_1^{y} - Q_2^{y} + Q_3^{y}\bigr). \label{eq:q4}
\end{align}
Следующая точка:
\begin{equation}
\label{eq:rk4}
\begin{aligned}
x_{n+1} &= x_n + \frac{1}{8}\bigl(Q_1^{x} + 3Q_2^{x} + 3Q_3^{x} + Q_4^{x}\bigr),\\
y_{n+1} &= y_n + \frac{1}{8}\bigl(Q_1^{y} + 3Q_2^{y} + 3Q_3^{y} + Q_4^{y}\bigr).
\end{aligned}
\end{equation}
В методичке в третьей стадии напечатано \(y_k - Q_1/2 + Q_2\). С этим коэффициентом погрешность на шаге падает только вдвое при уменьшении \(h\) вдвое, то есть метод становится первого порядка. У схемы с теми же узлами и весами \((1,3,3,1)/8\) третья стадия содержит \(-Q_1/3 + Q_2\). Именно эта формула имеет четвёртый порядок, и она использована в счёте.

На одном шаге, пока производные до пятого порядка ограничены,
\[
\|(x(t_n+h), y(t_n+h)) - (x_{n+1}, y_{n+1})\| = O(h^5),
\]
а накопленная погрешность к моменту \(T = Nh\) есть \(O(h^4)\).

Правая часть квадратична, поэтому вдали от особых точек решение за малое время уходит далеко. При счёте траектория обрывается, как только координаты перестают быть умеренными.

\section{Один шаг явно}

Начальные данные рядом с неустойчивым узлом, шаг из методички:
\[
x_0 = 1{,}2, \quad y_0 = -1, \quad h = 0{,}05.
\]
\[
f(1{,}2,-1) = 3{,}4^2 - 9 = 2{,}56,
\qquad
g(1{,}2,-1) = 9 - 3{,}2^2 = -1{,}24,
\]
\[
Q_1 = (0{,}128;\; -0{,}062).
\]
Вторая стадия в точке \((1{,}242667;\; -1{,}020667)\):
\[
Q_2 = (0{,}164602;\; -0{,}089233).
\]
Третья стадия в точке \((1{,}321935;\; -1{,}068566)\):
\[
Q_3 = (0{,}239109;\; -0{,}148257).
\]
Четвёртая стадия в точке \((1{,}402507;\; -1{,}121025)\):
\[
Q_4 = (0{,}320689;\; -0{,}214140).
\]
Формула~\eqref{eq:rk4} даёт
\begin{align*}
x_1 &= 1{,}2 + \frac{0{,}128 + 3\cdot 0{,}164602 + 3\cdot 0{,}239109 + 0{,}320689}{8}
     = 1{,}407478,\\[1mm]
y_1 &= -1 + \frac{-0{,}062 + 3\cdot(-0{,}089233) + 3\cdot(-0{,}148257) + (-0{,}214140)}{8}
     = -1{,}123576.
\end{align*}
За один шаг точка отошла от \((1,-1)\). Показатели узла положительны, \(18\) и \(6\), поэтому дальше уход ускоряется.

\section{Численные траектории}

Шаг \(h = 0{,}05\). В таблицах каждый второй узел, кроме короткого ухода от неустойчивого узла, где приведены все шаги до выхода из ограниченной области.

\subsection{Приближение к устойчивому узлу}

Старт \((-0{,}8;\; 1)\) лежит на расстоянии \(0{,}2\) от точки \((-1,1)\).

\begin{center}
\begin{tabular}{ccc}
\toprule
\(t\) & \(x\) & \(y\) \\
\midrule
\(0{,}0\) & \(-0{,}800000\) & \(1{,}000000\) \\
\(0{,}1\) & \(-0{,}925716\) & \(1{,}038651\) \\
\(0{,}2\) & \(-0{,}965752\) & \(1{,}028046\) \\
\(0{,}3\) & \(-0{,}982345\) & \(1{,}016564\) \\
\(0{,}4\) & \(-0{,}990509\) & \(1{,}009293\) \\
\(0{,}5\) & \(-0{,}994827\) & \(1{,}005136\) \\
\(0{,}6\) & \(-0{,}997167\) & \(1{,}002825\) \\
\(0{,}7\) & \(-0{,}998446\) & \(1{,}001552\) \\
\(0{,}8\) & \(-0{,}999148\) & \(1{,}000852\) \\
\(0{,}9\) & \(-0{,}999532\) & \(1{,}000468\) \\
\(1{,}0\) & \(-0{,}999743\) & \(1{,}000257\) \\
\bottomrule
\end{tabular}
\end{center}

Координата \(x\) идёт к \(-1\), \(y\) после небольшого выбега возвращается к \(1\). К моменту \(t = 1\) точка лежит в \(3\cdot 10^{-4}\)-окрестности равновесия.

\subsection{Уход от неустойчивого узла}

Старт \((1{,}2;\; -1)\), тот же шаг \(h = 0{,}05\).

\begin{center}
\begin{tabular}{ccc}
\toprule
\(t\) & \(x\) & \(y\) \\
\midrule
\(0{,}00\) & \(1{,}200000\) & \(-1{,}000000\) \\
\(0{,}05\) & \(1{,}407478\) & \(-1{,}123576\) \\
\(0{,}10\) & \(2{,}027366\) & \(-1{,}585611\) \\
\(0{,}15\) & \(5{,}894331\) & \(-4{,}865312\) \\
\bottomrule
\end{tabular}
\end{center}

К \(t = 0{,}2\) координаты уже порядка десятков тысяч: квадратичная правая часть разгоняет решение, покинувшее окрестность узла. Поэтому график окрестности строится с более мелким шагом и обрывается на расстоянии \(0{,}5\) от особой точки. Иначе одна ушедшая траектория растягивает оси, и локальная картина пропадает.

\subsection{Оценка погрешности по правилу Рунге}

На отрезке \([0;1]\) для начальных данных \((-0{,}8;\,1)\) решение посчитано с шагами \(h\) и \(h/2\). Для метода порядка \(p = 4\)
\begin{equation}
\label{eq:runge}
\varepsilon \approx \frac{z_h - z_{h/2}}{2^{p}-1} = \frac{z_h - z_{h/2}}{15}.
\end{equation}
При \(t = 1\):
\begin{align*}
h = 0{,}05:&\quad (-0{,}999743276;\; 1{,}000256698),\\
h = 0{,}025:&\quad (-0{,}999743406;\; 1{,}000256569).
\end{align*}
Разность координат имеет порядок \(1{,}3\cdot 10^{-7}\). Оценка~\eqref{eq:runge} даёт погрешность решения с шагом \(0{,}025\) около \(9\cdot 10^{-9}\) по каждой координате.

\section{Выводы}

\begin{enumerate}
  \item Система~\eqref{eq:system} имеет ровно четыре особые точки~\eqref{eq:points}. Они находятся решением \(2x-y = \pm 3\), \(x-2y = \pm 3\).
  \item Линеаризация~\eqref{eq:J} даёт:
  \((1,-1)\)~--- неустойчивый узел, показатели \(18\) и \(6\);
  \((-1,1)\)~--- устойчивый узел, показатели \(-6\) и \(-18\);
  \((3,3)\) и \((-3,-3)\)~--- сёдла, показатели \(\pm 6\sqrt{3}\).
  \item Счёт выполнен методом~2, формулы~\eqref{eq:q1}--\eqref{eq:rk4},
  шаг \(h = 0{,}05\).
  На траектории к \((-1,1)\) правило Рунге даёт погрешность порядка \(10^{-8}\).
  Траектория от \((1{,}2,-1)\) за время \(0{,}2\) уходит из окрестности узла.
  \item Единственная притягивающая точка~--- \((-1,1)\).
  Точка \((1,-1)\) отталкивает.
  Точки \((3,3)\) и \((-3,-3)\)~--- сёдла.
\end{enumerate}

\end{document}
